
\documentclass[aspectratio=169]{beamer}
\usetheme{Madrid}
\usecolortheme{seahorse}
\usepackage{booktabs}
\usepackage{hyperref}
\usepackage{listings}
\lstset{basicstyle=\ttfamily\small,breaklines=true}

% Increase default font sizes for a PPT-like look
\setbeamerfont{title}{size=\Huge,series=\bfseries}
\setbeamerfont{frametitle}{size=\LARGE,series=\bfseries}
\setbeamerfont{framesubtitle}{size=\large}

\title{Single Cycle Processor}
\subtitle{Course Project — CE222}
\author{Your Name}
\date{\today}

\begin{document}

\begin{frame}
  \titlepage
\end{frame}

\begin{frame}{Outline}
  \tableofcontents
\end{frame}

\section{Introduction}
\begin{frame}{Project Objectives}
  \begin{itemize}
    \item Implement a single-cycle RISC-like processor in Verilog
    \item Verify functionality with testbenches and waveform inspection
    \item Demonstrate datapath, control, and memory integration
    \item Present simulation results and example executions
  \end{itemize}
\end{frame}

\section{Architecture}
\begin{frame}{Top-level Block Diagram}
  % Replace 'block_diagram.png' with your generated diagram exported from drawing tool
  \centering
  \fbox{\parbox{0.85\linewidth}{\centering Missing \texttt{block\textunderscore diagram.png}\\Please add exported diagram here}}
  \vspace{0.5em}
  {\small (Insert annotated datapath: IF, ID, EX, MEM, WB)}
\end{frame}

\begin{frame}{Datapath Overview}
  \begin{itemize}
    \item Instruction Fetch: Program Counter, Instruction Memory
    \item Instruction Decode: Register File, Immediate Generator, Control Unit
    \item Execute: ALU, ALU Control, MUXes
    \item Memory Access: Data Memory (load/store)
    \item Write Back: Register update from ALU or Memory
  \end{itemize}
\end{frame}

\section{Modules}
\begin{frame}{Key Verilog Modules}
  \begin{itemize}
    \item \textbf{Instruction Memory:} \texttt{instruction\_memory.v}
    \item \textbf{Register File:} \texttt{register\_file.v}
    \item \textbf{Program Counter:} \texttt{pc.v}
    \item \textbf{ALU:} \texttt{alu.v}, \texttt{alu\_control.v}
    \item \textbf{Control Unit:} \texttt{control.v}
    \item \textbf{Data Memory:} \texttt{data\_memory.v}
  \end{itemize}
  \vspace{0.5em}
  % Files are located in the \texttt{Single\_Cycle\_Processor} folder of the repository.
\end{frame}

\begin{frame}{Register File}
  \begin{itemize}
    \item 32 registers, synchronous write, asynchronous read
    \item Ports: read1, read2, write, writeData, writeEnable
    \item Testbench: \texttt{register\_file\_tb.v} (use to validate read/write semantics)
  \end{itemize}
\end{frame}

\begin{frame}{ALU and ALU Control}
  \begin{itemize}
    \item ALU supports add, sub, and, or, slt, shifts, etc.
    \item \texttt{alu\_control.v} maps opcode/funct fields to ALU ops
    \item Unit-tested with \texttt{test\_alu} in repo
  \end{itemize}
\end{frame}

\begin{frame}{Control Unit}
  \begin{itemize}
    \item Decodes opcode to generate control signals (RegWrite, MemRead, MemWrite, ALUSrc, MemToReg)
    \item Single-cycle control: all signals valid in one clock cycle per instruction
    \item File: \texttt{control.v}
  \end{itemize}
\end{frame}

\section{Verification}
\begin{frame}[fragile]{Testbenches and Simulation}
  \begin{itemize}
    \item Use \texttt{iverilog}/\texttt{vvp} to simulate and \texttt{gtkwave} to view waveforms
    \item Primary testbench: \texttt{single\_cycle\_processor\_tb.v}
    \item Additional benches in repo: \texttt{instruction\_fetch\_tb.v}, \texttt{register\_file\_tb.v}, etc.
  \end{itemize}
  \vspace{0.5em}
  \begin{block}{Suggested Simulation Commands}
    \begin{verbatim}
# compile
iverilog -o simv *.v *_tb.v
# run
vvp simv
# view waveforms (if .vcd produced)
gtkwave single_cycle_processor.vcd
    \end{verbatim}
  \end{block}
\end{frame}

\begin{frame}{Example Waveforms}
  \centering
  % Replace 'waveform_example.png' with actual exported waveform images (PNG)
  \fbox{\parbox{0.9\linewidth}{\centering Missing \texttt{waveform\textunderscore example.png}\\Export VCD screenshots or PNG here}}
  \vspace{0.5em}
  {\small Show instruction execution, PC updates, register writes, memory accesses.}
\end{frame}

\section{Results}
\begin{frame}{Functional Results}
  \begin{itemize}
    \item Correct execution of sample instruction sequences verified via VCD
    \item Register read/write and memory load/store validated in testbenches
    \item Timing: single-cycle latency determined by longest combinational path (ALU + memory)
  \end{itemize}
\end{frame}

\begin{frame}{Performance / Observations}
  \begin{itemize}
    \item Pros: Simple control, easy to reason about, minimal pipeline hazards
    \item Cons: Single-cycle restricts clock frequency due to long-critical paths
    \item Potential improvements: pipelining, forwarding, hazard detection
  \end{itemize}
\end{frame}

\section{Conclusion}
\begin{frame}{Conclusion}
  \begin{itemize}
    \item Implemented a working single-cycle processor in Verilog
    \item Verified with testbenches and waveform inspection
    \item Next steps: add pipeline stages and hazard logic for performance
  \end{itemize}
\end{frame}

\begin{frame}{Appendix: Files of Interest}
  \begin{itemize}
    \item \texttt{single\textunderscore cycle\textunderscore processor.v} -- top-level (Single\textunderscore Cycle\textunderscore Processor folder)
    \item \texttt{single\textunderscore cycle\textunderscore processor\textunderscore tb.v} -- top-level testbench
    \item \texttt{instruction\textunderscore memory.v}, \texttt{data\textunderscore memory.v} -- memories
    \item \texttt{register\textunderscore file.v}, \texttt{alu.v}, \texttt{control.v}, \texttt{pc.v} -- core modules
  \end{itemize}
\end{frame}

\begin{frame}{How to Build the Presentation}
  \begin{itemize}
    \item Compile: \texttt{pdflatex presentation.tex} (run twice for TOC)
    \item Include exported PNGs for block diagram and waveforms in the same folder
    \item Recommended: use Google Slides / PowerPoint export if required
  \end{itemize}
\end{frame}

\end{document}
